\chapter{官能团是有机分子的“反应接口”}

\begin{kaichang}
去厨房和洗手台转一圈，你大概能找到这四样东西：一小瓶医用酒精，一瓶白醋，一瓶洗甲水（主要成分叫丙酮），还有一片阿司匹林。

先做个气味实验（只闻，别喝，也别混一起）：酒精有股熟悉的“医院味”，白醋酸得冲鼻子，洗甲水甜腻又刺鼻，阿司匹林片几乎没味道，但舔一下会酸。再看溶解性：酒精和水随便混，白醋更是泡在水里卖；洗甲水和水却分成两层，可是它一碰到泡沫塑料，塑料表面立刻发黏、塌陷。

这四样东西有一个共同的秘密：它们分子的骨架都是碳。上一章（第 36 章）刚说过，碳靠四只“手”搭出链、环和支链，搭出几百万种结构。可问题来了——既然骨架都是碳，为什么这几样东西的脾气差这么多？差别到底藏在哪里？

先猜一个答案。这一章结束时，你会发现：决定脾气的不是骨架，而是骨架上插着的几个小小的“接口”。
\end{kaichang}

\section{只看表面：差别到底有多大}

在深入分子之前，先把能看见、能测到的差别摆出来。四样东西都放在室温下：

医用酒精是无色液体，泼在手上凉得很快——它在迅速挥发，挥发要吸热（第 27 章的账）。它和水任意比例互溶，点得着，火焰是淡淡的蓝色。白醋也是无色液体， $5\%$ 左右的乙酸水溶液，倒在苏打粉上立刻冒泡，滴在大理石上会慢慢蚀出小坑。洗甲水挥发得比酒精还快，倒在桌上几十秒就只剩一层白雾似的凉气；它和水“谁也不理谁”，却能溶解泡沫塑料和指甲油。阿司匹林是白色药片，常温下安安静静，但在水里溶解得很勉强，泡半小时也只化开一小部分，化开的水尝着发酸。

注意一个共同点：这四样东西\textbf{都能燃烧}（药片烧起来会冒黑烟），烧完的气体通入石灰水都会变浑——都有二氧化碳生成。这说明它们分子里都有碳。可“都含碳”这个共同点，和“脾气差这么多”这个差别，摆在一起就成了本章真正的问题：同样的原料，是什么让成品如此不同？

\section{骨架是机身，接口才是按键}

把你家里的电器想一遍。手机、遥控器、游戏手柄，外壳形状千差万别，但你上手第一件事是找什么？找按键和接口——充电口、电源键、音量键。外壳决定它拿在手里什么感觉，按键和接口决定它\emph{能干什么}。

有机分子惊人地类似。第 36 章讲过，碳骨架可以由单键、双键、三键、环和支链搭出无穷花样。但对一个中等大小的有机分子来说，骨架上绝大部分的碳—碳键和碳—氢键性格相当“沉闷”：它们不爱反应，只是安安静静把分子撑起来。真正活跃、真正“管事”的，是骨架上某些特定位置的一小撮原子组合——氧和氢组成的一对、一个氧双键插在碳上、一个氮带着氢……化学家给这些组合起了名字：\textbf{官能团}。“官”是职能的“官”：官能团就是分子身上\textbf{反复出现、并且决定分子化学职能的部位}。

这话不是比喻，是可以验证的规律。把甲烷（\chem{CH_4}，天然气的主要成分）的一个氢换成一个“氧加氢”组合，得到乙醇（\chem{CH_3CH_2OH}，酒精）。骨架只多了一个碳、换了一个接口，结果：甲烷是难溶于水的可燃气体，乙醇却成了能与水任意混合、能消毒、能让人醉的液体。换掉一个接口，脾气整个改写。

再看一组更整齐的。骨架同样是两个碳，接口不同：

\begin{itemize}[itemsep=2pt]
\item 乙烷 \chem{CH_3CH_3}：没有任何特殊接口，气体，点得着，别的几乎什么都不会。
\item 乙醇 \chem{CH_3CH_2OH}：带羟基（\chem{-OH}），液体，溶于水，易燃。
\item 乙醛 \chem{CH_3CHO}：带醛基（碳氧双键加一个氢），刺鼻，是酒在体内“闯祸”的中间产物。
\item 乙酸 \chem{CH_3COOH}：带羧基（碳氧双键再加羟基），就是白醋里的醋酸，能腐蚀大理石，能让紫甘蓝汁变色。
\end{itemize}

四个分子，骨架几乎一样，性格天差地别。差别全部来自接口。这就是为什么有机化学不学“几百万种分子一个一个背”，而是学\textbf{十几种接口}——接口认得，分子的大致脾气就猜得到。

\section{接口凭什么这么横：电子分布不均}

官能团为什么偏偏爱反应？答案在第一册就已经备好了，我们把三样东西拼起来用：电负性（第 19 章）、分子间作用力（第 22 章）、能量与速率（第 27、29 章）。

回忆第 19 章：共价键里的共享电子不是平均分家，电负性大的原子把电子云往自己这边拉。碳和氢的电负性接近，所以碳氢键、碳碳键里的电子分布比较均匀，整条骨架像一段电压平稳的电线，不招惹谁。而氧和氮的电负性明显大于碳和氢。于是一旦骨架上插进氧或氮，电子云立刻被拉歪：氧那头电子偏多、带一点负，碳或氢那头电子偏少、带一点正。

这一歪，三件事跟着发生。

第一件，\textbf{分子间吸力变大}。羟基（\chem{-OH}）上，氧带负、氢带正，一个分子的氢正好能被邻居分子的氧吸引——这就是第 22 章讲过的氢键。分子之间手拉手，想变成气体就得先挣脱，沸点自然被抬高。

第二件，\textbf{溶解性跟着改写}。第 23 章说过，水分子自己也带正负极，还靠氢键结成网。带羟基、羧基的分子能插进这张网，所以酒精、醋酸和水随便混；而纯骨架分子（比如汽油里的烷烃）插不进去，就被挤出来分层——“相似相溶”说的其实就是“接口相似才相溶”。

第三件，也是最重要的，\textbf{反应有了明确的“着陆点”}。第 29 章讲过，反应先要有分子撞对地方。电子偏多的氧，天生吸引缺电子的伙伴；电子偏少的碳，天生吸引电子丰富的伙伴。于是别的分子一来，直奔这些带“正负标记”的接口——反应几乎总是在官能团上开场。这正是下一章（第 38 章）要讲“亲核”“亲电”的基础，这里先记住一句话：\textbf{极性给反应画出了靶心}。

\begin{qiaoliang}
接口对沸点的影响有多大？比一比两个分子量几乎相同的分子：丙烷 \chem{C_3H_8}（相对分子质量 44）沸点约 $-42\,^{\circ}\mathrm{C}$，乙醇 \chem{C_2H_5OH}（相对分子质量 46）沸点约 $78\,^{\circ}\mathrm{C}$。分量差不多，沸点差了 $120\,^{\circ}\mathrm{C}$！唯一的结构差别就是乙醇多了一个羟基。按分子间作用力的账来算，氢键给每个分子“加”的束缚，大约相当于要把分子从液体里拔出来需要多花一倍的能量——这就是 $120$ 度落差的来源。所以你在厨房里闻得到酒精味（它还在挥发），却永远烧不开一锅丙烷（它早跑光了）。
\end{qiaoliang}

\section{认接口，像认插头}

现在把最常用的接口摆成一张表。不用背，把它当成插头图鉴：下次见到结构式，对照着认。

\begin{table}[htbp]
\centering
\begin{tabular}{lll}
\toprule
官能团 & 结构特征 & 一句话性格 \\
\midrule
碳碳双键（烯） & \chem{C=C} & 电子扎堆的地方，容易被“加”上去 \\
碳碳三键（炔） & \chem{C\equiv C} & 比双键更拥挤、更活泼 \\
碳卤键 & \chem{C-X}（\chem{X} 为氟、氯等） & 碳端偏正，容易被顶替 \\
羟基（醇） & \chem{-OH} 接在碳骨架上 & 能成氢键，亲水，可被氧化 \\
醚键 & \chem{C-O-C} & 氧夹在中间，性格相对安静 \\
醛基 & \chem{-CHO} & 碳氧双键在链端，活泼易氧化 \\
羰基（酮） & \chem{C=O} 夹在碳链中间 & 与醛同宗，稍稳重 \\
羧基 & \chem{-COOH} & 羟基遇上羰基，氢容易被水“撬走”，显酸性 \\
酯基 & \chem{-COO-} & 羧酸的氢被换成碳链，常有果香 \\
氨基 & \chem{-NH_2} & 氮上电子多，显碱性，爱抓质子 \\
酰胺键 & \chem{-CONH-} & 蛋白质里连成肽链的那种接口 \\
\bottomrule
\end{tabular}
\end{table}

表里有两个家族值得多说一句。\textbf{醛、酮、羧酸、酯}其实是一家：它们共用“碳氧双键”这个核心（化学家叫它羰基），区别只在羰基两旁接的是氢、碳链、羟基还是“氧加碳链”。羰基的碳因为电子被氧拉走而偏正，是整个家族的“事故高发区”——第 40 章讲酒的代谢、果香和肥皂时，你会看到它们轮番登场。

另一个是\textbf{芳香环}。苯环（六个碳围成一圈、电子在整个环上流动）本身就是一个特殊的“接口家族”，它的电子分布和反应方式自成一派，复杂到要单独用一章来讲——第 41 章专门处理它，本章你只要知道：看到六元环里画着一个圆圈，那就是芳香环在打招呼。

\section{读图训练：五张“身份证”}

理论说够了，直接上手读开场那几样东西的结构简式。第 36 章教过你：结构简式把碳骨架摊开写，每个碳凑够四只手。现在多一步——\textbf{先把官能团圈出来}。

\textbf{乙醇}（医用酒精的主角）：\chem{CH_3CH_2OH}。两个碳的骨架，末尾一个羟基。就这一个接口，解释了它能与水混溶、能燃烧、还能被身体里的酶一步步氧化。

\textbf{乙酸}（白醋的酸味来源）：\chem{CH_3COOH}。骨架同样两个碳，接口换成了羧基。羧基里的羟基氢特别“坐不稳”——第 32 章讲过，酸就是容易给出质子的分子，乙酸到了水里，一部分分子会把 \chem{H^+} 交给水，所以白醋显酸性。

\textbf{丙酮}（洗甲水）：\chem{CH_3COCH_3}。三个碳，中间夹着羰基。羰基让它分子间有点吸力，却又没法像羟基那样互相拉成氢键大网，所以它沸点不高（约 $56\,^{\circ}\mathrm{C}$）、挥发飞快——这正是洗甲水“一擦就干、凉飕飕”的原因。同时它既有点极性又有碳骨架，水溶不了的油墨、塑料它却能溶开，所以去污力强。

\textbf{阿司匹林}（乙酰水杨酸）：一片药里同时插着两个接口——一个羧基（所以吞下去有点酸）、一个酯基。这两个接口不是随便选的：它的“原料”水杨酸本来带一个羟基，对胃刺激很大；化学家把羟基改造成酯基，刺激性降了下来，药效保留。改一个接口，就是一种药的诞生——第 41 章还会回到药物结构这个话题。

\textbf{咖啡因}（咖啡、茶、可乐里的那位）：骨架里两个含氮的环拼在一起，接口主要是环上的羰基和氮。氮的电子分布恰好让它能卡进大脑里某种“困意信号”的接收器，把信号挡在门外——你喝了咖啡不觉得困，不是咖啡因给了你能量，是它把“该困了”的通知拦截了。

看出门道了吗？读结构式的顺序永远是：\textbf{先找接口，再看骨架}。接口告诉你“它会干什么”，骨架告诉你“它多大、多重、多油”。

这套“认接口”的本事，价值远超化学考试。它正是现代医药和材料工业的日常工作语言。药物化学家改造一个候选药物时，最常做的事就是换接口：溶解度不够，在骨架上接一个羟基或把羧酸做成钠盐；在体内分解太快，把容易被酶攻击的酯基换成更稳的酰胺；想让分子穿过肠道壁，就减少亲水接口、增加碳骨架的比例。同一个有效成分，接口微调几处，吸收速度、副作用、保质期都可能整个改观。农药、染料、香料、洗衣液里的表面活性剂，走的也是同一条设计路线——先定骨架，再插接口，像给机身选装按键。

顺便说命名。有机化合物的名字看起来像咒语，其实是一套把接口写进名字的系统：后缀暴露接口（“醇”有羟基、“醛”有醛基、“酸”有羧基），前缀和数字告诉你骨架多长、接口插在哪。把它当\textbf{读图工具}用——从名字猜结构、从结构猜名字——就够了。没有谁要求你背出二十个碳的酸叫什么，正如没有地图软件要求你背下全国街道名。

\begin{shiyan}
\textbf{接口决定“跟谁混”——三种液体的溶解实验}

材料：三个透明杯子、水、食用油、医用酒精（或高度白酒）、几滴酱油、一支筷子。

步骤：三个杯子各装半杯水。第一杯滴几滴酱油，搅拌；第二杯加一勺食用油，搅拌后静置；第三杯加两勺酒精，搅拌。然后再往第二杯（油水混合物）里倒两勺酒精，搅拌，观察。

观察点：酱油（成分多带亲水接口）消失在水里；油（几乎纯碳骨架）浮在水面分层；酒精与水完全不分彼此；而倒进油水混合物后，酒精一部分进水、一部分进油，两层之间的界线会变模糊。

安全提示：全程常温、不点火、不品尝；酒精远离灶台。这个实验只改变“谁和谁混”，不产生任何新物质——你在亲眼验证“接口决定溶解性”。
\end{shiyan}

\section{“接口”这个想法是怎么被证实的}

\begin{zhengju}
官能团不是谁拍脑袋发明的分类法，它是一百多年争论和实验逼出来的概念。

19 世纪初，有机化学是一团乱麻：从动植物里分离出的物质成千上万，每个都起一个俗名，没人知道它们之间有什么关系。当时还流行“生命力论”——认为有机物只能由活的生物制造。1828 年，维勒在实验室里用无机物氰酸铵加热，得到了尿素——一种典型的“动物的”物质。生命力论开始裂缝：有机物并不神秘，它就是普通原子按规则搭出来的。

可规则是什么？1832 年，维勒和李比希研究苦杏仁油的一系列反应时发现：一组原子（他们叫“苯甲酰基”）在好几个反应里\textbf{作为一个整体}，从一种化合物搬到另一种化合物，像一块积木被拆下来又装上去。他们大胆提出：有机分子可以看成由若干“基”拼成，基的行为可以跨分子追踪。这个想法立刻遭到围攻——当时连原子是否真实存在都在争论，凭什么说分子内部还有“基”这种小团体？反对者说，基只是账本上的记号，不能当真。

争论的裁决来自两路证据。第一路是\textbf{同系物的规律性}：甲醇、乙醇、丙醇……每多一个 \chem{CH_2}，沸点、密度按几乎固定的台阶上升，而化学行为（能燃烧、能被氧化成酸）整个系列一模一样。如果性质由整个分子随机决定，不该出现这么整齐的台阶；台阶说明：骨架负责“量”的渐变，接口负责“质”的相同。第二路证据晚一些，但直接得多——\textbf{红外光谱}。20 世纪中期，化学家让红外光穿过样品，发现不同波长的光被有选择地吸收，而且吸收位置和官能团一一对应：凡是有羟基的分子，都在同一个波段附近“咬一口”光；凡是有羰基的，在另一个波段咬。原理是第 9 章的老朋友：分子里的键像弹簧，特定弹簧只吸收特定频率的光。羟基、羰基不管长在哪一个分子上，弹簧劲都差不多，所以“咬”的位置几乎不变。这等于给官能团拍了指纹照——它不再是有用的记账记号，而是分子内部真实存在、可以被仪器指认的局部结构。

今天，有机化学家拿到一个陌生分子的谱图，第一件事仍然是找这些“咬痕”——和你在这一章学的读图顺序一模一样：先找接口。
\end{zhengju}

\section{这个模型什么时候会失灵}

“认接口猜脾气”是有机化学最顺手的工具，但它有两个边界，你得心里有数。

第一，\textbf{同一个接口，在不同邻居身边，脾气会变}。羟基接在碳骨架上是醇（酒精），接在苯环上就成了苯酚——酸性比醇强得多，还有腐蚀性和毒性。羧酸里那个羟基，因为挨着羰基，氢就比醇的羟基容易电离得多。接口不是密封舱，骨架会通过化学键“遥控”接口的电子分布。精细的规律（取代基怎样推拉电子）是第 38、41 章的内容，这里记住定性结论就够了：接口决定大方向，邻居负责微调。

第二，\textbf{多个接口会协同，整体大于局部之和}。阿司匹林的两个接口互相影响它在胃酸里的溶解和分解；咖啡因的羰基和氮排成特定的空间阵列，才能恰好卡进受体。当分子大到一定程度——比如蛋白质——成百上千个接口在三维空间里折叠排布，会涌现出任何单个接口都没有的本事（催化、识别、运动）。那是本书后半程的主题，第三册见。

所以正确的心法是：\textbf{先用接口快速定位，再用整体结构修正判断}。只背接口表、不看分子全貌，就像只看按键就断言一台机器的全部性能。

\begin{wuqu}
“看到羟基（\chem{-OH}）就叫酒精。”——这是把接口当成了分子的全部。甲醇 \chem{CH_3OH} 和乙醇只差一个 \chem{CH_2}，羟基一模一样，但甲醇在体内会被氧化成甲酸，几毫升就能损伤视神经，工业酒精掺假致人失明的新闻都源于此。反过来，乙酸也有羟基（藏在羧基里），谁也不把白醋叫酒。羟基决定的是“能成氢键、可被氧化”这一类倾向，至于它把整个分子带向何方，还要看骨架和邻居。同理，“含羧基就是危险的强酸”也不成立——柠檬、苹果、酸奶里全是羧酸，它们的酸性温和（第 32 章讲过，强弱和浓稀是两回事）。接口是线索，不是判决书。
\end{wuqu}

\section{回到洗手台}

现在重新检查开场那四样东西。

医用酒精：骨架两个碳，一个羟基接口。羟基让它与水任意混溶，所以能配成消毒水；它还能钻进细菌体内捣乱（具体的破坏方式涉及蛋白质，第三册讲），这是消毒能力的来源。

白醋：同样两个碳，接口换成羧基，从“中性可燃液体”变成“能给出质子的酸”。酸冲鼻子的气味、能软化水垢的本事，都记在这一个接口头上。

洗甲水：三个碳夹一个羰基。羰基给了它恰到好处的“半极性”——足以溶开油墨和塑料，又不足以被水拉住，于是挥发飞快。

阿司匹林：一个羧基负责酸性和水溶性，一个酯基负责降低刺激、并在体内被水解回有效成分。两个接口分工合作，撑起一片药。

你最初的猜想对了几条？四条线索归结成一句话：\textbf{骨架是差不多的机身，接口才是按键}。这就是有机化学把几百万种分子压缩成一堂课的秘密。

再补一笔数量级的感受。一小杯约 \ud{15}{mL} 的纯酒精，密度约 \ud{0.79}{g/mL}，乙醇的摩尔质量是 \ud{46}{g/mol}，算下来约 $0.26$ 摩尔，也就是 $1.6\times 10^{23}$ 个分子——每个分子都带着一个羟基接口等着被处理。你肝脏里的酶大约每小时只能“接待”其中十分之一左右（约 \ud{10}{mL} 纯酒精），剩下的分子带着接口在血液里排队。所谓“喝多了”，从分子层面看，就是接口的处理队列排到了大脑门口。这也解释了为什么“解酒偏方”都不靠谱：第 29、30 章说过，反应速率和催化有其上限，酶的数量和速度不会因为喝茶、喝醋而改变。

\section{接口认全了，反应还远吗}

这一章你只学会了“认接口”，还没回答一个更硬的问题：接口上到底\emph{怎样}发生反应？为什么有的分子冲上去替换别人（取代），有的拆开双键往里加（加成）？电子是怎样一对一对搬家的？

好消息是，这些反应也不用一条条背。下一章（第 38 章）会给你一把万能钥匙：跟着电子走。弯箭头一画，谁推电子、谁拉电子、键在哪里断在哪里成，全部一目了然——亲核体扑向缺电子的接口，亲电体扑向电子扎堆的双键，有机反应瞬间从几千条背诵条目变成几种基本舞步。再往后，第 40 章带你用这把钥匙重新逛一遍醇、醛、酸、酯（酒的代谢、水果的香气、肥皂的去污都在那里），第 41 章则去拜访最特殊的一类接口——芳香环，看染料为什么有颜色、药物怎样被设计出来。

\begin{tiaozhan}
同样是羟基上的氢，为什么乙酸愿意把它交给水（显酸性），乙醇却几乎不肯？用本章的接口思想加第 32 章的酸碱定义可以推理大半。氢离开羟基后，氧上多出一对电子、带一个负电荷——这个负电荷能不能被“分摊”，决定了氢离开的难易。乙醇里，负电荷只能孤零零地压在一个氧上；乙酸的羧基里，负电荷可以在\textbf{两个氧之间}分摊（通过碳氧之间的电子重新排布），负担减半，所以乙酸的氢离开的倾向比乙醇大约强 $10^{11}$ 倍（用 p$K_a$ 衡量：乙酸约 $4.8$，乙醇约 $16$，相差 11 个数量级）。一个接口的邻居，把性质拉开了一百多亿倍——这就是“接口决定大方向、邻居负责微调”最极端的例证。跳过这段不影响主线。
\end{tiaozhan}

\section*{练一练}

\begin{sikaoti}
\item 下面三个分子的骨架都只有一到两个碳：\chem{CH_4}、\chem{CH_3OH}、\chem{HCOOH}。不看任何表格，只用本章的接口思想，按沸点从低到高给它们排序，并说明每个分子靠什么“接口”（或没有接口）分子间相互吸引。
\item 画出乙醛 \chem{CH_3CHO} 的结构简式，用 $+$ 号和 $-$ 号标出醛基上电子偏少和偏多的两端，再画一个水分子的氧端靠近它的示意图。图旁注明：$+$ 和 $-$ 表示电子云偏移，不是完整的离子电荷。
\item 己烷 \chem{C_6H_{14}}（汽油成分之一）和己醇 \chem{C_6H_{13}OH} 骨架一样长，只多了一个羟基。预测：哪一个更易溶于水？哪一个沸点更高？如果这个羟基从链端挪到链中间，你的预测会变吗？为什么？
\item 阿司匹林的“原料”水杨酸带羟基，成品把羟基改成了酯基。参照乙醇和乙醚（\chem{CH_3CH_2OCH_2CH_3}，羟基的氢被换成碳链）的性质差异，推理：这个改造会让分子形成氢键的能力变强还是变弱？这对药片在体内的行为可能有什么影响？
\item 画一张“分子体检表”：随便从家里找一种食品标签上的成分（如柠檬酸、苯甲酸钠、山梨酸钾），查它的结构式，圈出所有官能团，并给每个接口写一句“它可能带来的性质”。查不到完整结构式时，至少从名字里的“酸”“醇”“酯”后缀推断它含有什么接口。
\end{sikaoti}
